\section{Annealed closure of the full physical trace}
\label{sec:tpc43-annealed-trace}

For fixed \((\gamma,j,k)\), group the left coefficients by their
squarefree--kernel label:
\begin{equation}
 A^L_{\gamma,j,k}(s)
 :=\sum_{\alpha:s_{\alpha,j,k}=s}
 a^L_{\alpha,\gamma,j,k}\mu^2(n_{\alpha,j,k}).
 \label{eq:tpc43-grouped-left-coefficient}
\end{equation}
Equation \eqref{eq:tpc43-fused-Walsh-character} gives the exact Walsh
polynomial
\begin{equation}
 [\mathscr T_{k,\varepsilon}^L]_{\gamma,j}
 =\sum_s A^L_{\gamma,j,k}(s)\chi_s(\varepsilon).
 \label{eq:tpc43-coordinate-Walsh-polynomial}
\end{equation}

\begin{theorem}[Full trace at the annealed diagonal scale]
\label{thm:tpc43-annealed-trace}
For the complete physical alias family,
\begin{align}
 &\E_\varepsilon
 \sum_{k\in\cI_B}
 \bigl(
  \|\mathscr T_{k,\varepsilon}^L\|^2
  +\|\mathscr T_{k,\varepsilon}^R\|^2
 \bigr)
 \notag\\
 &\qquad\le \rho_X\cD_{\rm tr}
 \ll X^{o(1)}Q^2JK_B.
 \label{eq:tpc43-annealed-trace-bound}
\end{align}
The estimate retains the actual masks, orbit coordinates, source weights,
terminal cutoffs, and all aliases.  It does not use the projective layer
decomposition.
\end{theorem}

\begin{proof}
Walsh orthogonality gives
\begin{equation}
 \E_\varepsilon
 |[\mathscr T_{k,\varepsilon}^L]_{\gamma,j}|^2
 =\sum_s|A^L_{\gamma,j,k}(s)|^2.
 \label{eq:tpc43-one-coordinate-expectation}
\end{equation}
For one fiber, Cauchy--Schwarz and
\cref{lem:tpc43-kernel-fiber-bound} give
\begin{align}
 |A^L_{\gamma,j,k}(s)|^2
 &\le r_{j,k}(s)
 \sum_{\alpha:s_{\alpha,j,k}=s}
 |a^L_{\alpha,\gamma,j,k}|^2
 \mu^2(n_{\alpha,j,k})
 \notag\\
 &\le\rho_X
 \sum_{\alpha:s_{\alpha,j,k}=s}
 |a^L_{\alpha,\gamma,j,k}|^2
 \mu^2(n_{\alpha,j,k}).
 \label{eq:tpc43-fiber-Cauchy}
\end{align}
Sum over \(s,\gamma,j,k\) to obtain the left contribution
\(\rho_X\sum_k\cD_{k,L}\).  The same argument gives the right
contribution.  Now use \eqref{eq:tpc43-rho-bound} and the inherited
diagonal bound \eqref{eq:tpc43-inherited-trace-diagonal}.
\end{proof}

Within the annealed multiplicative model, the theorem removes the full
factor
\begin{equation}
 \frac{Q}{K_B}=X^{o(1)}J^2
 \label{eq:tpc43-annealed-closes-J2}
\end{equation}
that separates the deterministic overlap estimate from the sufficient
equality endpoint in TPC--42.  The statement is annealed: it concerns a
single globally multiplicative prime--sign field averaged over its prime
coordinates, not the distinguished all--minus point.

The same geometry acts directly on the minimal folded bank.  Let
\(\cK_0=\{0\}\), and for \(1\le a\le L_B\) let
\(\cK_a=\{a,a-(L_B+1)\}\).  Define
\begin{equation}
 [\mathscr U_{a,\varepsilon}^L]_{\gamma,j}
 :=\sum_\alpha\sum_{k\in\cK_a}
 a^L_{\alpha,\gamma,j,k}
 F_\varepsilon(d_\alpha)F_\varepsilon(n_{\alpha,j,k}),
 \label{eq:tpc43-random-folded-column}
\end{equation}
and define the right side symmetrically.

\begin{corollary}[Annealed closure of the actual minimal bank]
\label{cor:tpc43-annealed-folded}
One has
\begin{equation}
 \E_\varepsilon\sum_{a=0}^{L_B}
 \bigl(
  \|\mathscr U_{a,\varepsilon}^L\|^2
  +\|\mathscr U_{a,\varepsilon}^R\|^2
 \bigr)
 \le2\rho_X\cD_{\rm tr}
 \ll X^{o(1)}Q^2JK_B.
 \label{eq:tpc43-annealed-folded-bound}
\end{equation}
\end{corollary}

\begin{proof}
At fixed \((\gamma,j,a)\), expand the inner sum into atoms
\((\alpha,k)\) with \(k\in\cK_a\).  For each Walsh label there are at
most two alias choices, and for each alias there are at most \(\rho_X\)
rows.  Walsh orthogonality and fiberwise Cauchy--Schwarz therefore cost
at most \(2\rho_X\).  Summing the raw atomic coefficients over the
disjoint buckets gives exactly \(\cD_{\rm tr}\).
\end{proof}

\begin{corollary}[One--scale typical trace closure]
\label{cor:tpc43-one-scale-trace}
Let \(L_X\to\infty\) be any function.  Then
\begin{align}
 &\Pp_\varepsilon\left(
 \sum_k
 \bigl(
 \|\mathscr T_{k,\varepsilon}^L\|^2
 +\|\mathscr T_{k,\varepsilon}^R\|^2
 \bigr)
 >L_X\rho_X\cD_{\rm tr}
 \right)
 \le L_X^{-1}.
 \label{eq:tpc43-trace-Markov}
\end{align}
In particular, choosing any slowly growing \(L_X\to\infty\) with
\(L_X=X^{o(1)}\) yields the trace endpoint for all but an \(o(1)\)
proportion of prime--sign environments.
\end{corollary}

\begin{proof}
If \(\cD_{\rm tr}=0\), the annealed bound forces the trace energy to
vanish almost surely.  Otherwise apply Markov's inequality to that
nonnegative energy and use
\cref{thm:tpc43-annealed-trace}.
\end{proof}

\begin{remark}[A finite exact ensemble]
\label{rem:tpc43-finite-ensemble}
Only primes dividing the finitely many source and terminal integers enter
the packet.  Thus every expectation above is a normalized sum over a
finite cube.  No limiting probability space or independence between
different values of \(F_\varepsilon(n)\) is being assumed.
\end{remark}
